\chapter{Portafoglio di Markowitz}
\label{cap:markowitz}

\noindent\textbf{Classe:} QP convesso \hfill
\textbf{Script:} \texttt{python/lab06\_markowitz.py}

\section{Motivazione gestionale}

Quanto investire in ciascun titolo per bilanciare rendimento atteso e rischio? Il
modello di Markowitz è il QP più famoso della storia: minimizza la varianza del
portafoglio imponendo un rendimento minimo. Il risultato non è un numero ma una
\emph{frontiera efficiente}: il menu completo dei compromessi rischio-rendimento tra
cui il decisore sceglie in base alla propria avversione al rischio.



Il problema nasce in finanza (Markowitz, 1952, premio Nobel) ma la sua struttura
--- allocare un budget tra alternative rischiose e correlate --- riappare ovunque:
mix di progetti di R\&S, portafogli di clienti, perfino l'allocazione di capacità
produttiva tra mercati volatili. È anche il ponte naturale tra l'ottimizzazione e
la statistica: i dati del modello sono \emph{stime}, con tutte le conseguenze del
caso.

\textbf{In questo capitolo impareremo a:} (1) formulare e risolvere un QP convesso;
(2) costruire e leggere la frontiera efficiente; (3) capire l'effetto dei vincoli
di mandato; (4) diffidare delle stime dei rendimenti attesi.

\section{Costruzione guidata del modello}

\paragraph{Il problema a parole.} \emph{Decidiamo} la quota di capitale da investire
in ciascun titolo. \emph{L'obiettivo} è minimizzare il rischio del portafoglio,
misurato dalla varianza del suo rendimento. \emph{I vincoli}: investire tutto il
capitale (le quote sommano a 1), garantire un rendimento atteso minimo e rispettare
le quote minime e massime per titolo.

\subsection*{Insiemi e dati (input del modello)}
\begin{center}\small
\begin{tabular}{llp{7.2cm}}
\toprule
Simbolo & Tipo & Significato \\
\midrule
$I$ & insieme finito & titoli, con $n = |I|$ \\
$\boldsymbol\mu \in \R^n$ & vettore di reali & rendimenti attesi: $\mu_i$ per il titolo $i \in I$ \\
$\mat Q \in \R^{n \times n}$ & matrice, $\mat Q \succeq 0$ & covarianze dei rendimenti \\
$\bar r$ & $\in \R$ & rendimento minimo richiesto \\
$\ell_i,\, u_i$ & $\in [0, 1]$ & quote minima e massima per il titolo $i$ \\
\bottomrule
\end{tabular}
\end{center}

\subsection*{Variabili decisionali}
\begin{center}\small
\begin{tabular}{llp{7.2cm}}
\toprule
$\vet x \in \R^n$ & continue & $x_i$ = quota di capitale investita nel titolo $i \in I$ \\
\bottomrule
\end{tabular}
\end{center}

\begin{modello}[Minima varianza con rendimento minimo]
\begin{align}
\min\;& \sum_{i \in I}\sum_{j \in I} q_{ij}\, x_i\, x_j \notag\\
\text{s.t.}\;\; & \sum_{i \in I} \mu_i\, x_i \ge \bar r \label{eq:rendmin}\\
& \sum_{i \in I} x_i = 1, \qquad \ell_i \le x_i \le u_i \;\;\forall i \in I \label{eq:budget6}
\end{align}
Formulazioni equivalenti: $\max\, \sum_{i \in I} \mu_i x_i - \lambda \sum_{i \in I}\sum_{j \in I} q_{ij} x_i x_j$
(media-varianza) e $\max\, \sum_{i \in I} \mu_i x_i$ con varianza $\le \bar\sigma^2$
(massimo rendimento a rischio limitato).
Variando $\bar r$ (o $\lambda$, o $\bar\sigma$) si percorre la stessa frontiera.
\end{modello}

\begin{spiegazione}
\textbf{Obiettivo.} $\sum_{i \in I}\sum_{j \in I} q_{ij}\, x_i x_j$ è la varianza del rendimento
di portafoglio: contiene le covarianze, ed è qui che nasce la diversificazione ---
combinare titoli poco correlati abbassa il rischio complessivo sotto quello di ciascun
componente.

\textbf{Vincolo \eqref{eq:rendmin} (rendimento minimo).} Impone che il rendimento
atteso del portafoglio raggiunga la soglia $\bar r$: è la ``manopola'' con cui si
percorre la frontiera efficiente, e il suo moltiplicatore dice quanta varianza costa
ogni punto di rendimento in più.

\textbf{Vincolo \eqref{eq:budget6} (budget e limiti).} Con $x_i \ge 0$ la vendita allo scoperto è
vietata; il tetto $u_i$ è un tipico vincolo regolamentare o di mandato.

\textbf{Convessità.} Ogni matrice di covarianza è semidefinita positiva, quindi il QP
è convesso e l'ottimo trovato è \emph{globale certificato}.
\end{spiegazione}

\section{Esempio svolto a mano}

\begin{esempio}[Due titoli non correlati]
Volatilità $\sigma_1 = 20\%$, $\sigma_2 = 30\%$, correlazione nulla. Minimizzare la
varianza con $x_1 + x_2 = 1$ (senza vincolo di rendimento).

\textbf{Passo 1.} Sostituendo $x_2 = 1 - x_1$:
$f(x_1) = 0{,}04\,x_1^2 + 0{,}09\,(1 - x_1)^2$.

\textbf{Passo 2.} $f'(x_1) = 0{,}08\,x_1 - 0{,}18\,(1 - x_1) = 0
\;\Rightarrow\; x_1 = \dfrac{0{,}09}{0{,}04 + 0{,}09} = \dfrac{9}{13} \approx 69{,}2\%$.

\textbf{Passo 3.} Varianza del portafoglio:
$f(9/13) = 0{,}04 \cdot 0{,}479 + 0{,}09 \cdot 0{,}095 = 0{,}0277$, cioè volatilità
$16{,}6\%$: \textbf{meno rischioso di entrambi i titoli}. È l'effetto diversificazione
nella sua forma più pura: i pesi ottimi sono inversamente proporzionali alle varianze.
\end{esempio}

\section{Caso di studio}

Otto ETF settoriali (\emph{Exchange Traded Fund}: fondi quotati che replicano un
paniere di titoli di un settore), 60 rendimenti mensili simulati con un modello a un fattore di
mercato (file \texttt{dati/markowitz\_rendimenti.csv}); $\boldsymbol\mu$ e $\mat Q$ sono \emph{stimati}
da questi dati, come nella pratica. Statistiche annualizzate:

\begin{lstlisting}[style=output]
ENE: mu = 19,12%  vol = 18,65%      SAN: mu =  2,80%  vol = 11,93%
FIN: mu =  8,27%  vol = 19,59%      CON: mu = 11,25%  vol = 11,67%
TEC: mu =  1,58%  vol = 29,27%      UTL: mu = 11,99%  vol =  7,53%
IND: mu =  9,97%  vol = 16,60%      MAT: mu = 20,93%  vol = 23,93%
\end{lstlisting}

\begin{lstlisting}[style=python]
m = gp.Model("markowitz")
x = m.addVars(n, ub=u, name="x")          # quote, 0 <= x_i <= u
m.addConstr(x.sum() == 1, name="budget")
m.addConstr(gp.quicksum(mu[i]*x[i] for i in range(n)) >= r_min)
m.setObjective(gp.quicksum(Q[i, j]*x[i]*x[j]
                           for i in range(n) for j in range(n)),
               GRB.MINIMIZE)
m.optimize()
\end{lstlisting}

\section{Risultati}

\begin{lstlisting}[style=output]
Minima varianza globale : rendimento 11,06%, volatilita'  6,03%
Equipesato (1/n)        : rendimento 10,74%, volatilita' 10,13%
Composizione min varianza: ENE 5,2%  IND 1,9%  SAN 11,5%  CON 22,3%  UTL 58,3%
\end{lstlisting}

\begin{center}
\begin{tikzpicture}
\begin{axis}[lab, width=0.9\textwidth, height=7.4cm,
  title={Frontiera efficiente: la diversificazione domina i titoli singoli},
  xlabel={volatilità annua (\%)}, ylabel={rendimento atteso annuo (\%)}, xmin=4]
\addplot [teal, very thick] table [col sep=comma, x=vol, y=rend]
  {\figdat{cap06_front_libera}};
\addplot [arancio, thick, dashed] table [col sep=comma, x=vol, y=rend]
  {\figdat{cap06_front_limiti}};
\addplot [only marks, mark=*, mark size=1.8pt, grigiomedio,
  point meta=explicit symbolic, nodes near coords,
  every node near coord/.append style={font=\tiny, anchor=west, color=grigiomedio}]
  table [col sep=comma, x=vol, y=mu, meta=titolo] {\figdat{cap06_titoli}};
\addplot [only marks, mark=star, mark size=4pt, rossomattone]
  table [col sep=comma, x=vol, y=rend, discard if not={nome}{minvar}]
  {\figdat{cap06_speciali}};
\addplot [only marks, mark=diamond*, mark size=3pt, viola]
  table [col sep=comma, x=vol, y=rend, discard if not={nome}{equipesato}]
  {\figdat{cap06_speciali}};
\legend{frontiera senza limiti, frontiera con $u_i = 30\%$, titoli singoli,
        minima varianza, equipesato $1/n$}
\end{axis}
\end{tikzpicture}
\captionof{figure}{Piano rischio-rendimento: gli otto titoli (punti grigi), la
frontiera efficiente senza limiti (linea teal), la frontiera con tetto del 30\% per
titolo (tratteggiata arancione), il portafoglio di minima varianza (stella rossa) e
l'equipesato $1/n$ (rombo viola), che risulta dominato.}
\end{center}

\begin{center}
\begin{tikzpicture}
\begin{axis}[lab, width=0.9\textwidth, height=6.4cm, stack plots=y, grid=none,
  area legend,
  title={Composizione ottima lungo la frontiera},
  xlabel={rendimento richiesto $\bar r$ (\%)}, ylabel={quota (\%)},
  legend columns=8, legend style={at={(0.5,-0.28)}, anchor=north},
  ymin=0, ymax=100, enlargelimits=false]
\addplot [fill=teal!80, draw=none] table [col sep=comma, x=rend, y=ENE]
  {\figdat{cap06_composizione}} \closedcycle;
\addplot [fill=rossomattone!75, draw=none] table [col sep=comma, x=rend, y=FIN]
  {\figdat{cap06_composizione}} \closedcycle;
\addplot [fill=verde!75, draw=none] table [col sep=comma, x=rend, y=TEC]
  {\figdat{cap06_composizione}} \closedcycle;
\addplot [fill=arancio!80, draw=none] table [col sep=comma, x=rend, y=IND]
  {\figdat{cap06_composizione}} \closedcycle;
\addplot [fill=blunotte!75, draw=none] table [col sep=comma, x=rend, y=SAN]
  {\figdat{cap06_composizione}} \closedcycle;
\addplot [fill=grigiomedio!70, draw=none] table [col sep=comma, x=rend, y=CON]
  {\figdat{cap06_composizione}} \closedcycle;
\addplot [fill=viola!70, draw=none] table [col sep=comma, x=rend, y=UTL]
  {\figdat{cap06_composizione}} \closedcycle;
\addplot [fill=ocra!75, draw=none] table [col sep=comma, x=rend, y=MAT]
  {\figdat{cap06_composizione}} \closedcycle;
\legend{ENE, FIN, TEC, IND, SAN, CON, UTL, MAT}
\end{axis}
\end{tikzpicture}
\captionof{figure}{Composizione del portafoglio ottimo (aree impilate, in \%) al
crescere del rendimento richiesto $\bar r$: a sinistra domina il mix difensivo
(UTL, CON, SAN), verso destra il portafoglio si concentra progressivamente sui
titoli ad alto rendimento (MAT, ENE) fino al portafoglio monotitolo all'estremo.}
\end{center}

\begin{insight}
Tre messaggi dalla figura. (1)~Il portafoglio di minima varianza (volatilità 6\%) è
molto meno rischioso del miglior titolo singolo (UTL, 7,5\%) pur rendendo l'11\%.
(2)~L'equipesato $1/n$ --- la strategia ``non so nulla'' --- è nettamente dominato:
stessa area di rendimento ma quasi il doppio della volatilità.
(3)~Il tetto $u_i = 30\%$ taglia la parte alta della frontiera: il costo dei vincoli
di mandato si \emph{vede} come distanza tra le due curve.
\end{insight}

\section{Analisi di sensitività}

\begin{lstlisting}[style=output]
r_min =  6%: vol 6,03%   (vincolo NON attivo: coincide col min varianza)
r_min =  8%: vol 6,03%   (idem)
r_min = 10%: vol 6,03%   (idem)
r_min = 12%: vol 6,10%   d(varianza)/d(r_min) ~ 0,0221
\end{lstlisting}

\begin{spiegazione}[Un vincolo che non costa nulla]
Fino a $\bar r = 11\%$ il vincolo di rendimento è \emph{inattivo}: il portafoglio di
minima varianza rende già l'11,06\%, quindi chiedere ``almeno l'8\%'' non costa nulla
e il moltiplicatore è zero. Solo oltre l'11,06\% il vincolo morde e ogni punto di
rendimento in più si paga in varianza (moltiplicatore $\approx 0{,}022$ a
$\bar r = 12\%$). Riconoscere i vincoli inattivi evita di ``pagare'' vincoli gratis.
\end{spiegazione}

\begin{attenzione}[La fragilità delle stime]
Nei dati simulati il titolo TEC ha $\alpha$ vero dell'11\% annuo, ma su 60 mesi il
rendimento \emph{stimato} è 1,6\%: il rumore domina. È la lezione pratica più
importante di Markowitz: \textbf{le stime dei rendimenti attesi sono molto più
instabili di quelle delle covarianze}, e i portafogli ottimizzati inseguono gli errori
di stima. Rimedi usati in pratica: vincoli $u_i$, shrinkage delle stime, oppure
ottimizzare solo il rischio (minima varianza).
\end{attenzione}

\section{Esercizi}

\begin{esercizio}[verifica]
Verificare a mano l'esempio dei due titoli con $\rho = 0{,}5$: la formula diventa
$x_1 = (\sigma_2^2 - \rho\sigma_1\sigma_2)/(\sigma_1^2 + \sigma_2^2
- 2\rho\sigma_1\sigma_2)$. La diversificazione conviene ancora?
\end{esercizio}

\begin{esercizio}[frontiera]
Ricalcolare la frontiera con $u_i = 20\%$: quanto rendimento massimo si perde?
Rappresentare le tre frontiere ($u = 1$, $0{,}3$, $0{,}2$) nello stesso grafico.
\end{esercizio}

\begin{esercizio}[stabilità delle stime]
Rieseguire lo script con un seme casuale diverso (\texttt{default\_rng(1)}) e
confrontare le composizioni ottime a parità di $\bar r$: quanto cambiano? Quale
implicazione operativa ne traete?
\end{esercizio}

\begin{esercizio}[tracking error]
Il benchmark è l'equipesato $x^B_i = 1/8$. Minimizzare il tracking error
$\sum_{i \in I}\sum_{j \in I} q_{ij} (x_i - x^B_i)(x_j - x^B_j)$ con rendimento atteso $\ge 12\%$: quanto ci si deve scostare dal
benchmark?
\end{esercizio}

\begin{esercizio}[costi di transazione]
Partendo dal portafoglio equipesato, aggiungere il costo quadratico
$\kappa \sum_{i \in I} (x_i - 1/8)^2$ con $\kappa = 0{,}5$: come cambia la composizione
rispetto all'ottimo ``pulito''?
\end{esercizio}
